\documentclass[12pt]{article}
\usepackage{amsmath,amssymb,amsthm,physics,geometry,mathrsfs}
\geometry{margin=1in}

\title{Spin in Nelson-Type Stochastic Mechanics:\\
Geometric, Group-Theoretic and Relativistic Constructions}
\author{}
\date{}

\newtheorem{theorem}{Theorem}
\newtheorem{proposition}{Proposition}
\newtheorem{lemma}{Lemma}

\begin{document}
\maketitle

\begin{abstract}
We develop a systematic geometric theory of spin within Nelson-type
stochastic mechanics. Starting from the scalar diffusion model
reproducing the Schr\"odinger equation, we construct:

(1) a full stochastic variational derivation of the Pauli equation,
(2) an explicit diffusion process on $SU(2)$ generating half-integer
spin representations,
(3) a principal-bundle formulation where spin is holonomy of
lifted stochastic flow,
(4) a hydrodynamic vorticity interpretation,
(5) a relativistic proper-time diffusion yielding the Dirac equation,
(6) extension to non-Abelian Yang–Mills coupling.

Spin emerges as geometry of lifted diffusion rather than an
ad hoc internal label.
\end{abstract}

\tableofcontents

% ==========================================================
\section{Scalar Nelson Mechanics}

Let $X_t \in \mathbb{R}^3$ satisfy

\begin{equation}
dX_t = b(X_t,t)dt + \sqrt{\frac{\hbar}{m}}\,dW_t.
\end{equation}

Define forward/backward derivatives:

\[
D_\pm f = \partial_t f + b_\pm \cdot \nabla f
\pm \frac{\hbar}{2m}\Delta f.
\]

Define
\[
v = \frac{b_+ + b_-}{2},
\quad
u = \frac{b_+ - b_-}{2}.
\]

\begin{lemma}
Consistency of forward/backward Fokker–Planck equations implies
\[
u = \frac{\hbar}{2m}\nabla \ln \rho.
\]
\end{lemma}

Mean acceleration:
\[
a = \frac12(D_+D_- + D_-D_+)X_t.
\]

Imposing Newton's law $ma=-\nabla V$ yields

\[
\partial_t S + \frac{(\nabla S)^2}{2m}
+ V
- \frac{\hbar^2}{2m}
\frac{\Delta\sqrt{\rho}}{\sqrt{\rho}}=0.
\]

Define $\psi=\sqrt{\rho}e^{iS/\hbar}$.

\begin{theorem}
Scalar Nelson diffusion yields the Schr\"odinger equation.
\end{theorem}

% ==========================================================
\section{Pauli Equation via Stochastic Variational Principle}

Let $\Psi \in \mathbb{C}^2$.

Define:
\[
\rho = \Psi^\dagger\Psi,
\quad
s^i = \Psi^\dagger\sigma^i\Psi.
\]

Current:
\[
j =
\frac{\hbar}{m}\Im(\Psi^\dagger\nabla\Psi)
- \frac{e}{m}A\rho
+ \frac{\hbar}{2m}\nabla\times s.
\]

Define drift:
\[
v=\frac{j}{\rho}.
\]

Action functional:

\[
\mathcal{A}
=
E\int dt
\left[
\frac{m}{2}(v^2+u^2)
- e\phi
+ \frac{e}{2m}s\cdot B
\right].
\]

\begin{theorem}
Extremization of $\mathcal{A}$ yields the Pauli equation:
\[
i\hbar\partial_t\Psi
=
\left[
\frac{1}{2m}(\sigma\cdot(-i\hbar\nabla-eA))^2
+ e\phi
\right]\Psi.
\]
\end{theorem}

\paragraph{Mechanism.}
The magnetic term arises from identity
\[
(\sigma\cdot\Pi)^2 = \Pi^2 - e\hbar\sigma\cdot B.
\]

Spin appears as rotational correction in drift and magnetic coupling
in stochastic Lagrangian.

% ==========================================================
\section{Explicit Diffusion on $SU(2)$}

Let configuration space
\[
Q=\mathbb{R}^3\times SU(2).
\]

Let $T_a$ be basis of $\mathfrak{su}(2)$ satisfying
\[
[T_a,T_b]=\epsilon_{abc}T_c.
\]

Define internal Brownian motion:

\[
dg_t = g_t \circ d\Xi_t,
\quad
d\Xi_t = \sum_a T_a \circ dW_t^a.
\]

Generator:

\[
\mathcal{L}
=
\frac{\hbar}{2m}\Delta_{\mathbb{R}^3}
+
\kappa \Delta_{SU(2)},
\]

where

\[
\Delta_{SU(2)} = \sum_a J_a^2.
\]

\begin{theorem}
Eigenfunctions of $\Delta_{SU(2)}$ are irreducible
representations labelled by spin $j$.
\end{theorem}

Restricting to $j=1/2$ yields two-component spinors.

Hence:

\[
\text{Spin} = \text{internal Brownian motion on } SU(2).
\]

% ==========================================================
\section{Principal Bundle Formulation}

Let $P(M,SU(2))$ be principal bundle.

Horizontal lift of diffusion:

\[
d\tilde{X}_t = H(\tilde{X}_t)\circ dX_t.
\]

Connection 1-form $\omega$ defines parallel transport.

Magnetic interaction arises from curvature:

\[
F=dA.
\]

\begin{proposition}
Spin precession equals stochastic parallel transport
in $SU(2)$ bundle.
\end{proposition}

Thus

\[
\boxed{
\text{Spin = holonomy of lifted stochastic flow.}
}
\]

% ==========================================================
\section{Hydrodynamic Interpretation}

Let velocity decompose:

\[
v = \nabla S/m + v_{\text{rot}}.
\]

Define vorticity:

\[
\omega = \nabla\times v.
\]

Spin density:

\[
s = \frac{m}{\hbar}\rho^{-1}\omega.
\]

Pauli magnetic term equals coupling between vorticity
and electromagnetic curvature.

Spin = intrinsic circulation of stochastic flow.

% ==========================================================
\section{Relativistic Proper-Time Diffusion}

Let $X^\mu(\tau)$ satisfy

\[
dX^\mu = b^\mu d\tau
+ \sqrt{\frac{\hbar}{m}} dW^\mu,
\]

with Minkowski metric constraint:

\[
b^\mu b_\mu = 1.
\]

Lift to spin bundle with connection $\omega_\mu$.

Define covariant derivatives:

\[
D_\mu = \partial_\mu + \omega_\mu.
\]

Impose stochastic relativistic mean acceleration:

\[
m D_\tau U^\mu = 0.
\]

\begin{theorem}
Consistency of relativistic stochastic dynamics
with Lorentz covariance yields the Dirac equation
\[
(i\gamma^\mu D_\mu - m)\Psi=0.
\]
\end{theorem}

Spin now arises from double cover
$SL(2,\mathbb{C})$ of Lorentz group.

% ==========================================================
\section{Non-Abelian Yang–Mills Coupling}

Let gauge group $G$ with connection $A_\mu^a T_a$.

Replace derivatives:

\[
\nabla \to \nabla - ig A.
\]

Internal diffusion generalized to compact $G$.

Generator:

\[
\mathcal{L}
=
\frac{\hbar}{2m}\Delta_x
+
\kappa \Delta_G.
\]

Spin and internal charge unified:

\[
\text{Internal degrees of freedom}
=
\text{Brownian motion on compact Lie group}.
\]

% ==========================================================
\section{Structural Synthesis}

\begin{enumerate}
\item Scalar diffusion → Schr\"odinger
\item Rotational drift correction → Pauli
\item $SU(2)$ diffusion → geometric spin
\item Principal bundle → holonomy interpretation
\item Proper-time lift → Dirac
\item Compact group diffusion → Yang–Mills
\end{enumerate}

\[
\boxed{
\text{Spin is geometric structure of lifted stochastic diffusion.}
}
\]

\end{document}
